\section{One Bill Paid, One Left}
\label{sec:ch06-one-bill-paid-one-left}

\index{shareable@\code{"@shareable}!the package applies it unasked|(}%
Go back to the exported schema and read the paging type, which nothing in this
chapter touched. Its four descriptions are dropped here as everywhere else in
the book, and nothing else is:

\begin{graphqlsdl}
type PageInfo {
  hasNextPage: Boolean! @shareable
  hasPreviousPage: Boolean! @shareable
  startCursor: String @shareable
  endCursor: String @shareable
}
\end{graphqlsdl}

\index{shareable@\code{"@shareable}!on a type the library generates}%
Four fields, four directives, none of them written by anybody. \code{PageInfo}
is generated by the paging convention chapter~\ref{ch:data-without-n-plus-1}
turned on, and it is the same type in every subgraph that pages anything.
Chapter~\ref{ch:federation-model} said that two subgraphs declaring one field
fail composition unless both call it \code{@shareable}. The library knows that
this type can recur in other services that enable paging, so it declares
it shareable in advance.

\index{shareable@\code{"@shareable}!is a promise about values}%
That is the promise chapter~\ref{ch:federation-model} described being made on
your behalf. \code{PageInfo} is a value type carried by a particular
connection, not an entity fetched independently from another service.
Sharing its field definitions lets several services expose that connection
shape; it does not promise that different services return the same pages.

\index{node@\code{node}!did not get the same treatment}%
Now read \code{Query} in the same document. \code{node} and \code{nodes} are
there, exported by the global object identification
chapter~\ref{ch:schema-design} turned on, and neither carries
\code{@shareable}. Sessions and the future Speakers service enable this
feature and export both roots. Ratings will not. The collision comes from
those two opt-ins, not from a requirement that every subgraph expose nodes.

\index{composition!the collision this chapter does not meet}%
One collision the library saw coming and paid for. One it did not. The
difference is where the two features live: paging and federation are in the same
package family and \code{PageInfo} is generated by code that knows federation is
loaded, while \code{node} comes from a convention that has never heard of it.

\index{subgraph!a graph of one composes trivially}%
None of which costs anything yet. A collision needs two subgraphs, this chapter
has one. Its schema still has to pass composition's validation; it simply
has no second document to collide with yet. The bill falls due in
chapter~\ref{ch:composition}, which is the first chapter to run a composer, and
it is the concrete form of the debt chapter~\ref{ch:federation-model} recorded
against \code{Query.node} without paying.

\index{schema!read the exported file after every change}%
I would read the exported schema after switching this on even though nothing
made you. Two of the three things this section describes are directives the
package wrote into your public API without asking, and the third is one it
declined to write. None of them appears in your source, none of them is in a
compiler message, and all three change what a composer will say about you.
\index{shareable@\code{"@shareable}!the package applies it unasked|)}%
